\section{Dataset}

We use the Ship of Theseus Paraphrased Corpus \cite{tripto2024ship}, a controlled benchmark for studying authorial identity under iterative paraphrasing.\footnote{\url{https://github.com/tripto03/Ship_of_theseus_paraphrased_copus}}

\subsection{Corpus Structure}

The corpus contains seven topically diverse datasets: \textbf{CMV} (Reddit ChangeMyView arguments), \textbf{ELI5} (Reddit explanations), \textbf{SCI\_GEN} (scientific text), \textbf{TLDR} (Reddit summaries), \textbf{WP} (WritingPrompts fiction), \textbf{XSum} (BBC news summaries), and \textbf{Yelp} (business reviews). These datasets span a wide range of registers (informal Yelp reviews to formal scientific text), discourse types (argumentative CMV to narrative WP), document lengths (47 words average for TLDR to 512 for WP), and entity densities (5.9 entities per document for Yelp to 18.7 for XSum). This diversity enables testing whether decay patterns generalize across domains. Each dataset includes texts from seven source authors: one human and six LLMs (OpenAI, LLAMA, Tsinghua, BigScience, PaLM, Eleuther-AI).

Each text is paraphrased sequentially three times to produce a four-stage chain: the original text $T_0$, first paraphrase $T_1$, second paraphrase $T_2$, and third paraphrase $T_3$, where each $T_n$ is generated by paraphrasing $T_{n-1}$.

\subsection{Paraphrasing Models}

Seven paraphraser variants are used across four model families:

\begin{itemize}
    \item \textbf{ChatGPT}: Full-text paraphrasing via OpenAI's GPT model.
    \item \textbf{PaLM2}: Full-text paraphrasing via Google's PaLM2.
    \item \textbf{Dipper} \cite{krishna2023dipper}: A controlled paraphraser with three diversity settings: default, low, and high.
    \item \textbf{Pegasus} \cite{zhang2020pegasus}: A sentence-level paraphraser with two modes: full (all sentences paraphrased) and slight (25\% of sentences paraphrased).
\end{itemize}

\subsection{Preprocessing}

After loading all seven datasets, we applied light preprocessing: removing PaLM2 markdown formatting artifacts, normalizing unicode whitespace, and filtering texts with fewer than 10 words. We preserved casing, punctuation patterns, and sentence structure since these are the stylometric features under investigation in RQ1. After deduplication on (document key, version), the processed corpus contains \textbf{66,400 rows}: 3,020 originals at $T_0$ and approximately 21,000 rows each at $T_1$, $T_2$, and $T_3$. Table~\ref{tab:corpus_stats} summarizes the corpus.

\subsection{Experimental Splits}

The experimental design is split-aware across two axes to avoid leakage. Along the \emph{tier} axis, RQ2 classifiers are trained exclusively on $T_0$ features and evaluated on $T_1$--$T_3$, so the model never sees a paraphrased copy of a document it was trained on. Along the \emph{document} axis, RQ3 paraphraser-identification uses \texttt{GroupShuffleSplit} with an 80/20 train--test ratio keyed on document identity, so every paraphrase of a given source document is confined to a single split. This prevents the classifier from memorising a document's $T_0$ signature and recognising it in paraphrased form at test time. Classifier hyperparameters are selected via 5-fold stratified cross-validation on the $T_0$ training set only; the held-out 20\% document split and the $T_1$--$T_3$ tiers are touched only at final evaluation.

\begin{table}[h]
\centering
\small
\begin{tabular}{lrrr}
\toprule
\textbf{Dataset} & \textbf{Paraphrased} & \textbf{Train} & \textbf{Versions} \\
\midrule
CMV & 5,626 & 1,666 & 22 \\
ELI5 & 10,447 & 3,181 & 22 \\
SCI\_GEN & 10,360 & 2,780 & 22 \\
TLDR & 8,404 & 2,559 & 22 \\
WP & 10,289 & 2,754 & 22 \\
XSum & 10,472 & 3,243 & 22 \\
Yelp & 10,802 & 3,344 & 22 \\
\midrule
\textbf{Total} & \textbf{66,400} & \textbf{19,527} & -- \\
\bottomrule
\end{tabular}
\caption{Corpus summary after preprocessing. Each dataset has 22 version strings (1 original + 7 paraphrasers $\times$ 3 iterations).}
\label{tab:corpus_stats}
\end{table}